\section{Introduction}
\label{sec:intro}

The positions of the turbines in a wind farm are fixed at the planning stage and determine, for its lifetime, how much of the available wind energy it converts. An upstream turbine leaves behind a wake of slower and more turbulent air that reduces the power of the turbines downstream, so the layout sets the wake loss and thereby the gross energy yield, and shapes the collector system and grid connection~\cite{Srikakulapu2018,Hou2019}. As the wake loss is a share of the wake-free energy, a difference of a few tenths of a percentage point is, within a given wake model, a lifetime energy difference of that size. The wind farm layout optimization problem (WFLOP) places a given number of turbines in a site to minimize the wake loss subject to boundary and minimum-spacing constraints; we address only this wake-loss part of layout design, for farms of up to 36 turbines.

For swarm and evolutionary computation, the continuous WFLOP is also a representative constrained benchmark. Its $2N$ coordinates are continuous, the objective is multimodal and, with a top-hat wake, discontinuous; the constraints are a site boundary and $N(N-1)/2$ pairwise spacing constraints, so the feasible region is nonconvex and becomes a small part of the bounding box as $N$ grows; one evaluation is cheap enough for thousands of runs; and standard instances with published results exist, from the Kusiak--Song cases~\cite{Kusiak2010} to the IEA Wind Task~37 case studies~\cite{Baker2019}. These properties, which it shares with many engineering design problems, have made it a common test bed for new metaheuristics~\cite{Azlan2021}, including swarm methods such as particle swarm optimization (PSO)~\cite{Pookpunt2013,Asaah2021} and, more recently, metaphor-based algorithms such as the salp swarm algorithm (SSA)~\cite{Mirjalili2017,Rezk2019}. What such studies can show depends on how the comparison is run.

A new algorithm is usually justified by a comparison with baselines, and its claimed advantage is only as reliable as that comparison. Guidelines for experiments with metaheuristics~\cite{Hooker1995,Derrac2011,BartzBeielstein2020,LaTorre2021} and critiques of metaphor-based algorithms~\cite{Sorensen2015,Aranha2022} both point to the same weak spots: baselines run with poor or unreported settings, budgets that are not equal, constraint violations that are ignored in the ranking, and statistics that confuse the absence of a significant difference with equivalence. In the WFLOP, the GECCO layout competition~\cite{Wilson2018} and the IEA Wind Task~37 studies~\cite{Baker2019,Thomas2023} compared optimizers on a common objective, and reviews list many metaheuristics tested under differing budgets and settings~\cite{Azlan2021}; these studies compare methods as submitted or as published, not how the configuration of a baseline or the choice of control changes which method appears best (Section~\ref{sec:related}). Four issues matter here. First, the PSO setting $w=0.7$, $c_1=c_2=2$ (inertia weight and acceleration coefficients), which keeps the acceleration coefficients of the original PSO~\cite{Kennedy1995} and was used in a preliminary, unpublished comparison of ours, lies in the order-1 (mean) convergence region but outside the smaller order-2 (variance) region~\cite{Poli2009,Bonyadi2017,Cleghorn2018}. Second, budgets counted in iterations, or without the gradient evaluations of gradient-based solvers, are not equal. Third, infeasible layouts returned by penalty-based metaheuristics must be ranked consistently. Fourth, mean values without paired tests and multiplicity correction~\cite{Benavoli2016} may mislead, and a nonsignificant difference is not evidence of equivalence~\cite{Lakens2017}. For two-phase hybrids, a further question is whether the first phase contributes more than a random start would.

Our thesis is that the conclusions of such comparisons depend strongly on how the baselines are configured and how the comparison is controlled. Under a controlled protocol on the discretized Kusiak--Song benchmark at 6,030 evaluations from random starts, a simple two-phase design, PSO followed by the basic variable neighborhood search (VNS)~\cite{Mladenovic1997,Hansen2001} (PSO-VNS), and a well-configured PSO rank first and second of eight methods and differ little on average, whereas salp-swarm first phases give no gain larger than the 0.05~pp margin over random sampling in the farm (at the benchmark spacing $4D$ and at $5D$, not at $6D$; Section~\ref{sec:spacing}), and the Laplacian SSA (LX-SSA), previously proposed by Solanki and Deep~\cite{Solanki2023}, is worse as published. PSO-VNS is a sequential hybrid in the memetic sense~\cite{Neri2012}, not a new PSO variant.
% CHECK-FINAL [C01]: PSO-VNS best average rank; [C12]: constriction PSO \NPSOPos; CHECK-EXTRA [X39]/[X40]: difference inside +-m at seed and case level ("differ little on average"; no equivalence claim here)
% CHECK-FINAL [C53]/[C54]: SSA-VNS vs RSD-VNS no gain; CHECK-EXTRA [X53]; CHECK-FINAL [C55]: LX-SSA-VNS worse than RSD-VNS; [C17]: SSA-VNS lower than LX-SSA-VNS on the case means ("worse as published")
The methods and the individual protocol elements are established; the contribution is the controlled evidence showing which conclusions change when these choices are made explicit. Our contributions are:
\begin{enumerate}
\item \emph{Protocol.} Established practice~\cite{Derrac2011,BartzBeielstein2020,LaTorre2021} applied to the continuous WFLOP (equal budgets counting gradient evaluations, seed-paired runs and Holm-corrected tests), combined with three established elements whose joint use we have not found reported in WFLOP comparisons: feasibility-aware qualification and ranking (standard in constrained evolutionary computation~\cite{Deb2000,Wu2017}), equal-cost random-sampling controls for hybrids, and an equivalence analysis~\cite{Lakens2017,Benavoli2017} at the seed, case and cluster level with a post hoc margin of $\pm0.05$ percentage points (pp) of wake loss; applied to 68 benchmark cases, a Horns Rev~1 block and IEA37 Case Study~1 (Sections~\ref{sec:setup} and~\ref{sec:beyond}). The protocol is not specific to the WFLOP and carries over to other constrained continuous problems solved by penalty-based swarm and evolutionary methods.
% CHECK-EXTRA [X39]/[X41]: equivalence and minimal margins at seed, case and cluster level
\item \emph{Baseline configuration.} The PSO setting $w=0.7$, $c_1=c_2=2$ (run without velocity clamping) lies beyond the order-2 boundary (Proposition~\ref{prop:pso-stability}); in a sweep of $c_1=c_2$, feasibility starts to degrade at the boundary and falls progressively beyond it, and the setting reverses the order of PSO and the salp-swarm hybrids (Section~\ref{sec:psosetting}). A stability condition from PSO theory thus serves as a practical check of a baseline before a comparison is run.
% CHECK-THEORY [T01]/[T04]: order-2 violated (4 > 3.50); CHECK-CSWEEP [S02]-[S05]: onset in the step containing c*, accelerating ("falls progressively beyond it"); [S10]/[S11]: clamping |v| largely rescues (Section psosetting)
% CHECK-FINAL [C27]: tab:baseline pool: SSA-VNS and LX-SSA-VNS ahead of PSO with the old setting, PSO first with constriction
\item \emph{Component analysis.} Against random sampling in the farm (RSD-VNS), salp-swarm first phases give no gain larger than the 0.05~pp margin at the benchmark spacing $4D$ (and at $5D$; not at $6D$, Section~\ref{sec:spacing}), and LX-SSA as published (at twice the evaluations per iteration) is worse, whereas a PSO phase gives a clear gain; at 6,030 evaluations the VNS phase acts essentially as one compass-search descent (Section~\ref{sec:ablation}). A control that samples the bounding square instead of the farm is too weak: against it, the SSA phase shows a small gain that disappears against sampling in the farm, so the choice of control itself changes the verdict on a first phase.
% CHECK-FINAL [C53]/[C54]: SSA-VNS vs RSD-VNS no gain; CHECK-EXTRA [X53]: gain > m ruled out (seed, case); CHECK-FINAL [C55]: LX-SSA-VNS worse; CHECK-EXTRA [X34]; CHECK-FINAL [C56]/[C14]: PSO phase beats RSD-VNS and RS-VNS
% CHECK-DIAG [D08]-[D10], [D20]: first descent \NDDescentSharePSOVNS% of the Phase-2 gain; cycles \NDCycleEvalsPctPSOVNS% of Phase-2 evaluations for \NDCycleSharePSOVNS% of the gain
\item \emph{PSO-VNS versus PSO.} PSO-VNS has the best average rank on the $15^\circ$ benchmark, and its layouts keep it when re-evaluated (not re-optimized) with $5^\circ$ and $1^\circ$ direction bins, also under a Gaussian wake at $1^\circ$, but its average advantage over a well-configured PSO is small (equivalent at the post hoc margin $\pm0.05$~pp at the seed and case levels, not at the scenario-cluster level, with the $15^\circ$ rose; \ensuremath {-}0.043~pp, significant and not equivalent, after re-evaluation with a $1^\circ$ rose), and case-level differences go both ways. We report where it gains (feasibility on the Horns Rev~1 and Lillgrund blocks, larger Data Set~II layouts) and where it does not (feasible starts, IEA37), and derive recommendations (Sections~\ref{sec:overall}--\ref{sec:discussion}). Because the PSO phase of PSO-VNS reproduces a PSO run exactly for a given seed, this comparison isolates the contribution of the local-search phase.
% CHECK-FINAL [C01]; CHECK-DIR [F06]: best at 5 and 1 deg; [F09]: Gaussian at 1 deg PSO-VNS first (PSO first under Gaussian at 15 deg, C22, hence no "every" claim); [F10]: 1 deg \NFPairMeanOne pp, 90% CI \NFPairCIOne, p < 0.05; CHECK-EXTRA [X39]/[X42]: equivalent at seed and case level, not at cluster level (CR1, CR2, BM)
% CHECK-EXTRA [X43]: \NXHetBeyondPSOVNS and \NXHetBeyondPSO cases beyond +-m in each direction
% CHECK-EXTRA [X52]: Horns Rev feasibility, McNemar p = \NXHRMcNemarP; CHECK-FINAL [C32]/CHECK-EXTRA [X19]: Data Set II N >= 10: \NCmPSOVNSvsPSOdsIILargeWins of \NCmPSOVNSvsPSOdsIILargeN lower, p_Holm <= \NXHolmPSOVNSvsPSOdsIILargeP; CHECK-DIR [F14]: survives 1 deg
% CHECK-FINAL [C28]: feasible starts: \NFbFeasBestRank first, PSO-VNS second; [C40]: best published IEA37 layouts better (strict); CHECK-DIR [F44]: also with projected published layouts
\end{enumerate}

\emph{Organization of the paper.} Section~\ref{sec:related} reviews related work. Section~\ref{sec:model} states the wake, power and constraint model; Section~\ref{sec:prelim} describes PSO and its stability conditions, the SSA, LX-SSA and the basic VNS; and Section~\ref{sec:hybrid} the two-phase template, PSO-VNS and the random-sampling controls. Section~\ref{sec:setup} gives the evaluation protocol. Section~\ref{sec:results} reports the effect of the baseline setting and the overall comparison, Section~\ref{sec:ablation} the component analysis, Section~\ref{sec:beyond} tests beyond the benchmark (larger budgets, feasible starts, Horns Rev~1, Lillgrund, IEA37) and Section~\ref{sec:robustness} the sensitivity to the modeling assumptions. Section~\ref{sec:discussion} discusses the findings and gives recommendations for benchmarking, Section~\ref{sec:limits} discusses threats to validity and Section~\ref{sec:conclusion} concludes.
